% GENERATED FILE. Edit canonical lessons, then run npm run generate:books.
% canonical-source-hash: fnv1a64:50df2fd5b4325476
% canonical-lessons: TE-C395-kanuka, TE-C395-anduceta, TE-C395-drstya, TE-C395-anaga, TE-C395-pratipadika

\chapter{Giving Reasons}
\label{ch:te-giving-reasons}
\hlchaptermodality{395}

\begin{chapteropening}
\textbf{By the end of this chapter:} \emph{I can say so, therefore and in view of.}

Complete the last lesson of chapter 395: I can say so, therefore and in view of.
\end{chapteropening}

\section[\texorpdfstring{\te{కనుక} (kanuka)}{కనుక (kanuka)}]{\te{కనుక} (kanuka) — so; hence}

\label{lesson:TE-C395-kanuka}



\begin{quote}
Before the new one: say the Telugu for to justify, then the Telugu for exact.
\end{quote}

\subsection*{You'll want to know: \te{కనుక}}
\textbf{\te{కనుక}} — \emph{kanuka} — \textquotedblleft{}so; hence\textquotedblright{}.

\begin{quote}
\te{వర్షం} \te{పడింది}, \te{కనుక} \te{మేము} \te{బయటకు} \te{వెళ్ళలేదు}.
\end{quote}

\begin{quote}
\emph{varṣaṁ paḍindi, kanuka mēmu bayaṭaku veḷḷalēdu.}
\end{quote}

\begin{quote}
It rained, so we did not go out.
\end{quote}

It opens the result, after the reason, as \te{అందుకే} does.

\begin{cousinweb}
\textbf{\te{కనుక}} (\emph{kanuka}) is a Telugu connective of reason whose origin is uncertain; DEDR gives it no article of its own.
\end{cousinweb}

\begin{grammarlens}[title={What you've built}]
One more word for giving reasons and opinions.
\end{grammarlens}

\subsection*{Your turn}
\begin{itemize}
\raggedright
  \item \emph{Say it:} \emph{kanuka}
  \item \emph{Say it:} \emph{kanuka}, once more
  \item \emph{Say it:} \emph{varṣaṁ paḍindi, kanuka mēmu bayaṭaku veḷḷalēdu.} — It rained, so we did not go out.
  \item \emph{Recall:} say the Telugu for to justify, then the Telugu for exact, then say \emph{kanuka} again
\end{itemize}

\begin{tcolorbox}[breakable,colback=teal!4,colframe=teal!35!black,title={Before you move on}]
What does \te{కనుక} mean? (\textquotedblleft{}So; hence\textquotedblright{}.) Say it once more.
\end{tcolorbox}
\section[\texorpdfstring{\te{అందుచేత} (anducēta)}{అందుచేత (anducēta)}]{\te{అందుచేత} (anducēta) — for that reason; therefore}

\label{lesson:TE-C395-anduceta}



\begin{quote}
Before the new one: say the Telugu for exact, then the Telugu for so.
\end{quote}

\subsection*{You'll want to know: \te{అందుచేత}}
\textbf{\te{అందుచేత}} — \emph{anducēta} — \textquotedblleft{}for that reason; therefore\textquotedblright{}.

\begin{quote}
\te{అతనికి} \te{జ్వరం} \te{వచ్చింది}, \te{అందుచేత} \te{పనికి} \te{రాలేదు}.
\end{quote}

\begin{quote}
\emph{ataniki jvaraṁ vaccindi, anducēta paniki rālēdu.}
\end{quote}

\begin{quote}
He had a fever, and for that reason he did not come to work.
\end{quote}

\te{అందు}, in that, + \te{చేత}, by: by that. A cousin of \te{అందువల్ల}, which you know.

\begin{cousinweb}
\textbf{\te{అందుచేత}} (\emph{anducēta}) is native Telugu: \emph{andu}, \textquotedblleft{}in that, there\textquotedblright{}, built on the far-pointing base \emph{a-} (DEDR 1), and \emph{cēta}, \textquotedblleft{}by\textquotedblright{}, the ending that marks the means.
\end{cousinweb}

\begin{grammarlens}[title={What you've built}]
One more word for giving reasons and opinions.
\end{grammarlens}

\subsection*{Your turn}
\begin{itemize}
\raggedright
  \item \emph{Say it:} \emph{anducēta}
  \item \emph{Say it:} \emph{anducēta}, once more
  \item \emph{Say it:} \emph{ataniki jvaraṁ vaccindi, anducēta paniki rālēdu.} — He had a fever, and for that reason he did not come to work.
  \item \emph{Recall:} say the Telugu for exact, then the Telugu for so, then say \emph{anducēta} again
\end{itemize}

\begin{tcolorbox}[breakable,colback=teal!4,colframe=teal!35!black,title={Before you move on}]
What does \te{అందుచేత} mean? (\textquotedblleft{}For that reason; therefore\textquotedblright{}.) Say it once more.
\end{tcolorbox}
\section[\texorpdfstring{\te{దృష్ట్యా} (dṛṣṭyā)}{దృష్ట్యా (dṛṣṭyā)}]{\te{దృష్ట్యా} (dṛṣṭyā) — in view of; considering}

\label{lesson:TE-C395-drstya}



\begin{quote}
Before the new one: say the Telugu for so, then the Telugu for for that reason.
\end{quote}

\subsection*{You'll want to know: \te{దృష్ట్యా}}
\textbf{\te{దృష్ట్యా}} — \emph{dṛṣṭyā} — \textquotedblleft{}in view of; considering\textquotedblright{}.

\begin{quote}
\te{వర్షం} \te{దృష్ట్యా} \te{పాఠశాలకు} \te{సెలవు} \te{ఇచ్చారు}.
\end{quote}

\begin{quote}
\emph{varṣaṁ dṛṣṭyā pāṭhaśālaku selavu iccāru.}
\end{quote}

\begin{quote}
In view of the rain, they gave the school a holiday.
\end{quote}

Formal; it follows a noun: \te{వర్షం} \te{దృష్ట్యా}, in view of the rain.

\begin{cousinweb}
\textbf{\te{దృష్ట్యా}} (\emph{dṛṣṭyā}) is borrowed whole from Sanskrit \emph{dṛṣṭyā}, \textquotedblleft{}by the sight (of)\textquotedblright{}, the instrumental case of \emph{dṛṣṭi}, \textquotedblleft{}seeing, sight\textquotedblright{} (Monier-Williams).
\end{cousinweb}

\begin{grammarlens}[title={What you've built}]
One more word for giving reasons and opinions.
\end{grammarlens}

\subsection*{Your turn}
\begin{itemize}
\raggedright
  \item \emph{Say it:} \emph{dṛṣṭyā}
  \item \emph{Say it:} \emph{dṛṣṭyā}, once more
  \item \emph{Say it:} \emph{varṣaṁ dṛṣṭyā pāṭhaśālaku selavu iccāru.} — In view of the rain, they gave the school a holiday.
  \item \emph{Recall:} say the Telugu for so, then the Telugu for for that reason, then say \emph{dṛṣṭyā} again
\end{itemize}

\begin{tcolorbox}[breakable,colback=teal!4,colframe=teal!35!black,title={Before you move on}]
What does \te{దృష్ట్యా} mean? (\textquotedblleft{}In view of; considering\textquotedblright{}.) Say it once more.
\end{tcolorbox}
\section[\texorpdfstring{\te{అనగా} (anagā)}{అనగా (anagā)}]{\te{అనగా} (anagā) — that is to say; namely}

\label{lesson:TE-C395-anaga}



\begin{quote}
Before the new one: say the Telugu for for that reason, then the Telugu for in view of.
\end{quote}

\subsection*{You'll want to know: \te{అనగా}}
\textbf{\te{అనగా}} — \emph{anagā} — \textquotedblleft{}that is to say; namely\textquotedblright{}.

\begin{quote}
\te{రేపు} \te{అనగా} \te{శుక్రవారం} \te{మాకు} \te{సెలవు}.
\end{quote}

\begin{quote}
\emph{rēpu anagā śukravāraṁ māku selavu.}
\end{quote}

\begin{quote}
Tomorrow, that is Friday, is a holiday for us.
\end{quote}

A bookish cousin of \te{అంటే}; it explains the word before it.

\begin{cousinweb}
\textbf{\te{అనగా}} (\emph{anagā}) is native Dravidian, a form of \emph{anu}, \textquotedblleft{}to say\textquotedblright{} (DEDR 868), the verb behind \te{అని} and \te{అంటే}.
\end{cousinweb}

\begin{grammarlens}[title={What you've built}]
One more word for giving reasons and opinions.
\end{grammarlens}

\subsection*{Your turn}
\begin{itemize}
\raggedright
  \item \emph{Say it:} \emph{anagā}
  \item \emph{Say it:} \emph{anagā}, once more
  \item \emph{Say it:} \emph{rēpu anagā śukravāraṁ māku selavu.} — Tomorrow, that is Friday, is a holiday for us.
  \item \emph{Recall:} say the Telugu for for that reason, then the Telugu for in view of, then say \emph{anagā} again
\end{itemize}

\begin{tcolorbox}[breakable,colback=teal!4,colframe=teal!35!black,title={Before you move on}]
What does \te{అనగా} mean? (\textquotedblleft{}That is to say; namely\textquotedblright{}.) Say it once more.
\end{tcolorbox}
\section[\texorpdfstring{\te{ప్రాతిపదిక} (prātipadika)}{ప్రాతిపదిక (prātipadika)}]{\te{ప్రాతిపదిక} (prātipadika) — a footing; a basis}

\label{lesson:TE-C395-pratipadika}



\begin{quote}
Before the new one: say the Telugu for in view of, then the Telugu for that is to say.
\end{quote}

\subsection*{You'll want to know: \te{ప్రాతిపదిక}}
\textbf{\te{ప్రాతిపదిక}} — \emph{prātipadika} — \textquotedblleft{}a footing; a basis\textquotedblright{}.

\begin{quote}
\te{మార్కుల} \te{ప్రాతిపదికన} \te{బహుమతులు} \te{ఇస్తారు}.
\end{quote}

\begin{quote}
\emph{mārkula prātipadikana bahumatulu istāru.}
\end{quote}

\begin{quote}
Prizes are given on the basis of marks.
\end{quote}

\te{ప్రాతిపదికన}, on the basis of, follows a noun in the genitive.

\begin{cousinweb}
\textbf{\te{ప్రాతిపదిక}} (\emph{prātipadika}) is borrowed from Sanskrit \emph{prātipadika}, a grammar word for \textquotedblleft{}the crude form or base of a noun\textquotedblright{} (Monier-Williams); Telugu widened it to any basis or footing.
\end{cousinweb}

\begin{grammarlens}[title={What you've built}]
One more word for giving reasons and opinions.
\end{grammarlens}

\subsection*{Your turn}
\begin{itemize}
\raggedright
  \item \emph{Say it:} \emph{prātipadika}
  \item \emph{Say it:} \emph{prātipadika}, once more
  \item \emph{Say it:} \emph{mārkula prātipadikana bahumatulu istāru.} — Prizes are given on the basis of marks.
  \item \emph{Recall:} say the Telugu for in view of, then the Telugu for that is to say, then say \emph{prātipadika} again
\end{itemize}

\begin{tcolorbox}[breakable,colback=teal!4,colframe=teal!35!black,title={Before you move on}]
What does \te{ప్రాతిపదిక} mean? (\textquotedblleft{}A footing; a basis\textquotedblright{}.) Say it once more.
\end{tcolorbox}
